%% SN Computer Science manuscript skeleton -- STAGE 1 (foundation only).
%%
%% This file is an initialized skeleton, not a draft. Section bodies contain
%% short scope notes, not prose. See ../MANUSCRIPT_PLAN.md for the outline
%% this skeleton follows, ../EVIDENCE_MAP.md for the claim-to-evidence table,
%% and ../README.md for compilation instructions.
%%
%% Template: Springer Nature sn-jnl.cls v3.1 (December 2024), vendored
%% unmodified in ../template/. Reference style: sn-basic (Numbered), matching
%% SN Computer Science's numeric-citation submission guidelines.
%% Peer review model: single-blind (author identity included; SN Computer
%% Science guidelines confirm single-blind, unlike the double-blind JDIQ_2026
%% submission this manuscript partially descends from).

\documentclass[pdflatex,sn-basic]{sn-jnl}

\usepackage{graphicx}
\usepackage{amsmath,amssymb,amsfonts}
\usepackage{booktabs}
\usepackage{multirow}
\usepackage{algorithm}
\usepackage{algorithmicx}
\usepackage{algpseudocode}

\raggedbottom

\begin{document}

%% ---------------------------------------------------------------------
%% Title: three candidates under consideration -- see MANUSCRIPT_PLAN.md
%% Section 1 for the full rationale. Working title used below is T1.
%% ---------------------------------------------------------------------
\title[Structural Consistency Is Not Retrieval Utility]{Structural
  Consistency Is Not Retrieval Utility: An Exact-and-Heuristic Audit of
  Preference-Graph Repair for Multi-Ranker Retrieval}

\author*[1]{\fnm{Soroush} \sur{Vahidi}}\email{sv96@njit.edu}

\affil*[1]{\orgdiv{Ying Wu College of Computing},
  \orgname{New Jersey Institute of Technology},
  \orgaddress{\city{Newark}, \state{NJ}, \country{USA}}}

%% ---------------------------------------------------------------------
%% Structured abstract (Purpose / Methods / Results / Conclusion), per
%% SN Computer Science submission guidelines: 150-250 words total, no
%% subheadings rendered as separate sections (bold run-in labels only),
%% no citations or equations inside the abstract. PLACEHOLDER -- not
%% drafted this stage; see MANUSCRIPT_PLAN.md Section 2 for the intended
%% content of each part.
%% ---------------------------------------------------------------------
\abstract{\textbf{Purpose:} [Stage 2 drafting. State the research question:
whether enforcing acyclicity in derived multi-ranker preference graphs --
via heuristic or exact minimum-weight feedback-arc-set repair -- yields a
statistically supported improvement in retrieval effectiveness, and what
that implies about using structural consistency as a proxy for retrieval
utility.]

\textbf{Methods:} [Stage 2 drafting. Summarize the audited pipeline
(construction, repair, extraction, evaluation), the four-benchmark
canonical package, the three vote-construction regimes, greedy vs.\ exact
SCIP repair, and the Holm-corrected paired-inference protocol.]

\textbf{Results:} [Stage 2 drafting. State the central null/conditional
result with the headline Holm-corrected cell counts and the exact-repair
confirmation, without overstating effect direction.]

\textbf{Conclusion:} [Stage 2 drafting. State the restrained conclusion:
structural repair is real and measurable, but is not a reliable surrogate
for downstream retrieval utility; implications for practitioners and
future preference-graph studies.]}

\keywords{preference graphs, feedback arc set, minimum-weight feedback arc
  set, graph repair, rank aggregation, information retrieval, retrieval
  evaluation, data quality}

\maketitle

%% =======================================================================
%% SECTION SKELETON
%% Each section below is a scope note for Stage 2+ drafting, not prose.
%% See MANUSCRIPT_PLAN.md Section 4 (section-by-section outline) for the
%% full rationale, and EVIDENCE_MAP.md for exact evidence citations.
%% =======================================================================

\section{Introduction}
\label{sec:introduction}
%% Scope: motivate preference graphs as derived data artifacts; state the
%% precise research question (see MANUSCRIPT_PLAN.md Section 2); state the
%% central contribution list (<=4 contributions, MANUSCRIPT_PLAN.md
%% Section 3); explicitly disclaim a state-of-the-art reranking framing.

\section{Related Work}
\label{sec:related_work}
%% Scope: rank aggregation / feedback-arc-set literature; probabilistic
%% pairwise models (Bradley-Terry, Rank Centrality); graph-free fusion
%% (RRF, CombSUM, Borda); LLM-as-judge / pairwise-ranking positional-bias
%% literature; data-quality-for-derived-artifacts framing. Reuse candidate:
%% IJCS draft Related Work + JDIQ Background and Related Work (see
%% MANUSCRIPT_PLAN.md Section 8, reuse plan).

\section{Preference-Graph Construction and Repair}
\label{sec:method}
%% Scope: formalize G_q=(V_q,E_q,w_q); the three vote-construction regimes
%% (ms2 / ms1 / ms1_drop_mutual); the MWFAS repair objective (greedy AND
%% exact SCIP, presented as co-equal, not exact-as-footnote per the
%% Iran JCS C7/C8 lesson); ranking extraction (Copeland, balance,
%% graph-free baselines). Reuse candidate: IJCS draft Method section
%% (graph formalism, Algorithm 1/2, regime table) reconciled with JDIQ's
%% seven-dimension audit-taxonomy framing.

\section{Experimental Design}
\label{sec:experimental_setup}
%% Scope: four benchmarks (SciDocs, FiQA, HotpotQA, BRIGHT); stored
%% BM25/TF-IDF/MiniLM scores; canonical vs. larger-pool (P>k) evaluation
%% cells; Holm-corrected paired nDCG inference as the primary confirmatory
%% test; bootstrap / exact-repair / alternative-pool / power analysis as
%% robustness checks, not separate headline families. Bound the six-query
%% real-LLM pilot explicitly as a directional addendum, not a second
%% confirmatory study (see MANUSCRIPT_PLAN.md Section 6, forbidden claims).

\section{Results}
\label{sec:results}
%% Scope subsections (see MANUSCRIPT_PLAN.md Section 5 for exact
%% table/figure list and evidence paths):
%%   5.1 Construction dominates graph structure (BM25 raw-vs-normalized
%%       edge-weight share; cyclicity by regime).
%%   5.2 Repair is structurally active but retrieval gains do not survive
%%       Holm correction (canonical + larger-pool families).
%%   5.3 Exact SCIP repair confirms the null is not a greedy-heuristic
%%       artifact (headline reviewer-facing result).
%%   5.4 Robustness: alternative pools, added baselines, power/MDE,
%%       narrow equivalence results.

\section{Discussion}
\label{sec:discussion}
%% Scope: state the restrained thesis explicitly (see Section 7 below);
%% separate structural consistency from downstream retrieval utility as
%% the paper's organizing distinction; practical audit-logic guidance for
%% practitioners; brief, explicitly bounded real-LLM addendum discussion.

\section{Limitations}
\label{sec:limitations}
%% Scope: finite power for small effects; benchmark/ranker scope; bounded
%% real-LLM evidence; candidate-pool sensitivity; no claim of transfer to
%% learned fusion, cross-encoder reranking, or online evaluation.

\section{Conclusion}
\label{sec:conclusion}
%% Scope: restate the central thesis in one paragraph; state what future
%% work would be required to overturn it (not what is planned).

\section{Data Availability and Reproducibility}
\label{sec:data-availability}
%% Scope: public-benchmark citations; repository URL; reproduction-script
%% pointers; explicit statement that stored intermediates and fixed
%% query-ID lists back every numerical claim (no paid API calls required
%% to regenerate reported statistics).

\backmatter

\section*{Declarations}

\subsection*{Funding}
[Stage 2: confirm funding statement -- likely ``Not applicable'' per the
IJCS draft precedent; verify before submission.]

\subsection*{Conflict of interest}
The author declares no conflict of interest.

\subsection*{Data and code availability}
The datasets used in this study are publicly available from their
original sources and are cited in the manuscript. The code and processed
artifacts used to generate the reported results are available at
\url{https://github.com/SoroushVahidi/consistency-aware-llm-rankin}.

\bibliographystyle{sn-basic}
\bibliography{references}

\end{document}
